% ch-notverified.tex -- verification: the explicit list of what has not been checked
\chapter{検証していないこと}
\label{ch:notverified}

確認済みの事項だけを並べたマニュアルを読むと、読者はそれ以外もすべて確認済みだと思い込みやすい。この種の思い込みは代償が大きい。誤りは黙って起こり、文書を信頼する人ほど陥り、結果が誤っていると判明して初めて表に出る。そこで本章では、確認して\textbf{いない}ことをはっきり述べる。以下の項目はいずれも欠陥ではなく、未検証の量が正しいことも十分ありうる。ただし、どの項目についてもYuzenは証拠を示していない。答えに頼る前に、利用者自身が証拠を用意すべき箇所である。検証に向けた最初の一歩が見えている場合は、それも挙げる。

\section{物理とモデル}
\label{sec:notverified:physics}

\begin{itemize}
  \item \textbf{2次元（および軸対称）の接触線。}\cref{sec:verification:disjoining,sec:verification:spreading}の内容はすべて1次元である。2次元の分離圧カーネルは、次元に依存しない本体を1次元のものと共有している。しかし、平衡液滴による検証も濡れ広がりの測定も、2次元では行っていない。最初の一歩は、\cref{sec:verification:disjoining}の構成を2次元で繰り返すことである。つまり、\textbf{解析的な軸対称平衡形状}の上で残差速度を測り、格子を細かくすると2次で減少することを確かめる。
  \item \textbf{マランゴニ流れと基板の凹凸は製造解でしか検証していない。}すべての項を有効にすると、どちらも2次で収束する（\cref{sec:verification:mms}）。ただし、どちらも閉じた形の解とは比較していない。前者については、表面張力勾配のもとにある平坦な膜の線形安定性が手がかりになると考えられる。\val{"CZ"}の製造解では濃度が面内だけでなく厚さ方向にも変化する。そのため、表面値（表面張力、マランゴニ応力）、分布$\mu(c(\zeta))$、写像による鉛直移流は、これらの値が互いに異なる状況で確認できている。対象外は蒸発である。製造解の分布は蒸発の表面条件を満たさない。
  \item \textbf{\val{"CZ"}の検証相手は鉛直一様モデルであり、解析解ではない。}濃度が膜厚に影響しえない極限では、\val{"CZ"}と\val{"WM"}は独立に書いた経路で同じ問題を解く。両者は1次元でも2次元でも相対値で約$10^{-6}$まで一致し、体積は厳密に保存される（\cref{sec:disc:positivity-concz}。テストは\code{concz\_vs\_wm}と\code{concz\_vs\_wm\_2d}）。整合性の確認としては強力だが、既知の答えに対する検証とは別物である。$x$--$y$対称性、$H_0$と$L_0$の選び方に対する不変性、1次元形と2次元形の一致も確認している（\cref{sec:verification:equivalence}）。ただし、これらで確かめられるのは離散化の構造であって、精度ではない。厚さ方向の拡散は、上面で一定流束を受ける平板の級数解と比較しており、誤差は$N_z$について2次で減少する（\cref{sec:verification:analytic}）。移流は鉛直一様モデルと比較している。乾燥する平坦な膜では、蒸発、$D(c)$による拡散、$\zeta$写像の鉛直移流がともに分布を形作る。この分布は独立な解と比較している（\code{dry\_cz}）。確認していないのは、平坦でない膜の上での同様の分布である。そこでは面内の流れと写像が同時に作用する。
  \item \textbf{状態ファイルのないリスタートは近似である。}状態ファイルを伴わない場ファイル（たとえば手で書いたもの）から始めると、適応時間刻み、\file{History.dat}の累計、定常状態のカウント、\val{"CZ"}の修正子が使う前ステップの濃度は、いずれも初期状態からやり直しになる。このため、このように継続した\file{History.dat}の\code{EvaporatedVol}、\code{OutflowedVol}、\code{VolErr}は、$t = 0$ではなくリスタート時点を基準とする（\cref{sec:ref-numerics:restart}）。
  \item \textbf{接触角ヒステリシスは扱わない。}動的接触角がCox--Voinov則から外れる領域も扱わない。
  \item \textbf{\val{"CZ"}にはナビエすべりを実装しておらず}、その入力は受け付けない。\val{"CZ"}は膜厚方向の速度分布を解くので、すべりを扱うにはそこに一様なすべり速度を加える必要がある。鉛直一様モデルのすべりは、毛管減衰率$\sigma(h_0^{3}/3 + \lambda h_0^{2})k^{4}/\mu$、すべての項を含む製造解、平衡液滴の静止状態、閉じた箱によって検証している（\cref{sec:verification:analytic,sec:verification:mms}）。確認していないのは、すべりを伴う濡れ広がり則（\cref{sec:verification:spreading}はすべりなし）と2次元の接触線である。すべりがあると、乾いた側のセル界面移動度は1次でしか消えない（\cref{sec:disc:positivity-slip}）。
\end{itemize}

\section{数値計算}
\label{sec:notverified:numerics}

\begin{itemize}
  \item \textbf{時間精度次数はすべての項を有効にして測ったが、スイッチの組み合わせをすべて試したわけではない。}\cref{sec:verification:mms}の製造解は、各モデルのすべての項を同時に含めた状態で$\Delta t$について2次である。ほかの項と組み合わせたときだけ2次になる項があっても、これでは気づかない。マランゴニ項がない場合、\val{"CZ"}の修正子が使う濃度は解いて求めるのではなく、直前の2ステップから外挿する。2次元の鉛直一様モデルの濃度は、1次元の拡散係数の10分の1で確認している。風上化の散逸が拡散誤差を打ち消すセルペクレ数を避けるためである。その代わり、\code{ctest}は線形問題で2次元濃度の漸近的な2次精度を確認する。平坦な膜の上で斜めの正弦波を移流・拡散させる問題である（\code{oblique\_wm}、\cref{sec:verification:schemes}）。1次元のパラメータですべての項を有効にした製造解が2次に近づく過程は、\code{ctest}に含まれていない。
  \item \textbf{制限付き対流スキームが有界になるのはクーラン数1以下に限られる}。クランク・ニコルソン法もこの条件を必要とする。ソルバーはこれを超えたステップを報告するが、刻みを縮めはしない。時間刻み制御もこの条件を考慮しない（\cref{sec:ref-numerics:convection}）。同梱のケースでは、膜厚の式でこれを超えるのが1つある。\file{nozzle\_1d\_rotating}（4.6）で、膜が壁際に盛り上がり、体積力と毛管圧が釣り合う箇所である。ステップが100倍長い\file{stress\_all}は0.94に達し、その他は0.03未満にとどまる。Korenのリミッターには折れ点があり、そこでは膜厚の式のヤコビアンが片側的になる。このリミッターを用いたときのニュートン法の収束は、\cref{sec:verification:schemes}のケースでしか確認していない。同梱のケースには本物の濃度前線をもつものがない。そのため、制限付きの\key{Concentration}スキームの利点は\file{tools/advection.py}でしか示していない。また\val{"CZ"}では、こうしたスキームは膜厚方向の分布の極値を切り落とす。この極値は流束ゼロ条件によって生じるものである（\cref{tab:schemes-mms}）。
  \item \textbf{\val{"CZ"}のマランゴニ結合の収束性は、平坦な膜についてしか導出していない。}膜と濃度は、中点の濃度が落ち着くまで、Aitken緩和を用いて中点で解き直す（\cref{sec:time:newton}）。濃度が成層した平坦な膜については、1回のパスのゲインを導出しており、測定値とも\SI{10}{\percent}以内で一致する。ゲインが1を超えるのは、モデルの波が$2/\Delta t$より速く成長する場合に限られる（\cref{sec:verification:czpasses}）。粘度が濃度に依存すると、膜のすべての層をつなぐ結合が加わる。この結合のゲインは導出しておらず、1つの製造状態で測っただけである。その状態では、状態を約1ラジアン動かすステップでパスが落ち着かなくなる。パスが八回以内に落ち着かないステップは失敗とし、刻みを縮める。固定刻みの場合は実行を止める。測定した例以外の長い実行でこれがどれほど起こるかは明らかでない。$\sigma_c\,\partial c/\partial\zeta > 0$で成層し、拡散係数の小さい膜は、モデル上、波長数マイクロメートルの波に対して不安定であり、これを打ち切るのは拡散だけである（\cref{sec:verification:schemes}）。成層した製造解は、格子で解像できるすべての波長（2倍細かい格子も含む）が安定になる拡散係数で作っている。したがって、こうした膜については何も明らかでない。この種の膜を計算すると、数セルの尺度に現れるのはこの不安定性を格子とスキームが打ち切った結果であり、収束した答えがどうなるかを示す確認はない。
    また、適応時間刻みを駆動する打ち切り誤差の推定値が実際の誤差に対応するかどうかも、\code{adaptive\_dt}（\cref{sec:verification:features}）以外では検証していない。
  \item \textbf{\val{"amgcl"}バックエンドはテストスイートの対象外である。}\val{"eigen"}との一致も、1スレッドと2スレッドの差の大きさも測定している（\cref{sec:verification:backends,sec:verification:amgclthreads}）。ただし、これは手作業で実行した単一の回帰ケースであり、\code{ctest}がコミットのたびに確かめる性質ではない。AMGCLビルドでも\code{ctest}は通る。このとき\code{restart\_roundtrip}の2次元の実行では\val{"auto"}が\val{"amgcl"}に解決されるが、確かめているのはリスタートであってバックエンドではない。\code{equivalence}は\val{"eigen"}を明示的に指定している。\cref{sec:verification:linalg}では、スレッド数を変えてもビット単位で一致することを\emph{テストで保証している}と述べた。この主張は\val{"eigen"}に限ったものである。\val{"amgcl"}は構造上ビット単位で再現しないので、この保証の対象にはなりえない。
  \item \textbf{\val{"FreshInlet"}と\val{"Outlet"}は互いに照らして確認しただけで、既知の流入とは比較していない。}流入口が課す値と各辺での振る舞いは、鏡映または転置した実行、および1次元の実行と照合している（\cref{sec:verification:bcio,tab:equivalence}）。流入口・流出口を通る体積と溶質の収支は、ステップがそこを通したセル界面流束と丸め誤差の範囲で一致することを確認している（\cref{sec:verification:analytic}）。流入分布が解析的に得られる問題との比較はない。流入口は、初期場のうちセル界面の両側にある2点の平均を保持する。初期分布が線形でない箇所では、この値は3次補間と$O(\Delta x^2)$だけ異なる。\val{"CZ"}は\val{"FreshInlet"}を受け付けない。
\end{itemize}

\section{テストスイートで検出できないこと}
\label{sec:notverified:mutation}

ここまでの確認は、コードが正しいかどうかを問うものであった。本節では、コードが正しくなかった場合に確認がそれを検出できるかを問う。そのために、物理のソースに単一の変更を108個、1つずつ加え、それぞれについて\code{ctest}を実行した。変更の種類は、符号の反転、添字を1つずらす、2分の1の係数を変える、\code{+=}を\code{=}にする、ループの境界をずらす、である。当初検出できたのは52個であった。残る56個が明らかにした穴を塞ぐため、ヤコビアンのテスト、すべての項を有効にした製造解、ノズル・蒸発の上限・流入口と流出口・乾燥分布の確認、および\file{tools/features.py}を書いた。これらはそのうち14個を検出する（うち三つはその後書き直したコードにあった）。さらに、流入口の値のテスト（\cref{sec:verification:bcio}）と、面内で濃度が変化する\val{"CZ"}の製造解が六個を検出する。膜厚と基板を含まない\file{.vts}からのリスタート、バイナリリスタートのマーカー探索、蒸発上限のしきい値、\val{"negPow"}の縁の確認がさらに四個を検出し、合計24個となる。流入口に対する五個の変更のうち四個は隅のゴースト値しか変えず、鉛直一様モデルのどの式もこの値を読まない。マーカー探索はどのリスタートからも観測できない。そのため、これらを捉えるのは単体テストだけである。残る32個に、テストされていない振る舞いは一つもない。

\begin{itemize}
  \item 13個は、ループや添字を配列の末尾の1つ先へ動かすものである。リリースビルドはもっともらしい値のまま走り切ってしまう。これを検出するのはサニタイザビルドであり、CIで1つ実行している（\cref{sec:verification:sanitizers}）。
  \item 11個は、テストしうるものを何も変えない。誰も読まないゴースト点を余分に訪れる（または飛ばす）ようになったループ、決して通らないフォールバック、リリースビルドでは除かれる\code{assert}、どのモデルも呼ばないコードである。
  \item 2個は誤差の大きさだけを変え、次数は変えない。1つのセル界面での3次補間の風上側を変えるもので、対象は\val{"CZ"}の鉛直移流と2次元の遠心力項である。どちらの場合もスキームは2次のままである。見分けるには、1つの解像度で参照解と数パーセントの精度で比較するしかない。
  \item 6個はメッセージ、カウンタ、出力にある。初期ステップが最大値を超えるという注意、ニュートン反復の総数、ピン留め確認の報告、定常状態の期間、場ファイルの隅の速度、使われていない変換である。
\end{itemize}

最後にカバーした四個から、欠陥が1つ見つかった。\key{NegPowRmax}より外側で\val{"negPow"}分布のブレンドが一定に保たれず、ブレンド幅の四倍の周期で振動し、縁での値の0.75まで下がっていた（\cref{sec:verification:analytic}）。

これらのテストはこの標本に合わせて書いたものなので、このスコアからは、まだ見ていない欠陥についてほとんど何も判断できない。そこで、これらのテストを入れた時点のソースから、同じ種類の変更で第二の標本を抽出した。最初の108個はいずれも含めていない。内訳は、モデルと数値計算のソース、リスタート、場のメタデータ、対流の設定にある3\,344通りの変更から一様無作為に100個（シード20260930）、入出力コードの残りにある312通りから25個（シード20261001）である。この標本については、実行し終えるまでテストを一切書かなかった。抽出ツールは文字列リテラルやエラーメッセージを含む行を飛ばすので、キーの解析の大部分はどちらの標本にも含まれない。

\begin{itemize}
  \item 最初の100個のうち、\code{ctest}は68個を検出した。残りのうち19個はどの結果も変えようがない（内部の判定の後に置かれたゴースト点の判定、どのモデルも呼ばないコード、\code{assert}、厳密に相殺する項）。したがって検出率は81個中68個（\SI{84}{\percent}）であり、CIのサニタイザビルドを含めると81個中72個（\SI{89}{\percent}）となる。2つの層の大きさで重み付けすると、それぞれ\SI{80}{\percent}と\SI{85}{\percent}である。
  \item 出力コードからの25個のうち、\code{ctest}は7個を検出した。サニタイザを含めると22個中10個である。見逃したのは結果ファイルに関するものであった。\file{.vts}のヘッダのオフセットの誤り（ParaViewはこれに従うが、リスタートは従わない）、誤った位置から書き出した速度や点の高さ、\val{"CZ"}の場ファイルの\code{Uz}列、製造解の範囲チェックである。
\end{itemize}

その後、次の確認を追加した。ヘッダを通して読んだ\file{.vts}を同じ時刻の場ファイルと比較するもの（\code{restart\_roundtrip}）、範囲チェック（\code{input\_guards}）、\val{"CZ"}でのプローブ位置（\code{features}）、Korenスキームを用いた成層\val{"CZ"}膜（\code{mms}）、強い回転のもとでの2次元の対称性（\code{equivalence}）である。これらは生き残っていたミュータントのうち20個を検出し、第一の標本は81個中73個（サニタイザを含めて77個）、第二の標本は22個中22個となった。この過程で欠陥が1つ見つかった。1次元の鉛直一様モデルの実行で、プローブファイルに記された点が、値を取った点より1セル手前を指していたのである。第一の標本に残るのは、テストの分解能より小さい変更が1つ（線形ソルバーの下限でニュートン法が1反復早く止まる）、出力ファイルのゴーストの隅またはゴースト列しか変えないものが三つ、サニタイザが検出する範囲外アクセスが四つである。

第三の標本は、コードの多くを書き直した後に抽出した。書き直したのは、風上化した流束の壁面セル界面、保存形の\val{"CZ"}濃度とその壁、時計、溶媒の限界、入力のより厳格な読み取り、リスタート状態である。変更の種類は同じで、それ以前に抽出した234個はいずれも含めていない。内訳は、第二の標本の最初の100個が属するファイル（今回はリスタート状態も含む）にある3\,815通りの変更から100個（シード20261004）、入出力コードの残りの309通りから25個（シード20261005）である。それぞれについて\code{ctest}全体を実行し、生き残ったものはCIのサニタイザビルドでもう一度実行した。

\begin{itemize}
  \item 100個のうち、\code{ctest}は74個を検出した。うち四個は制限時間の超過で、十個（すべて製造解の生成された項にある）は\code{mms\_source\_current}で検出した。17個は結果を変えようがない（どの式も読まないゴースト値、\code{assert}、厳密に等しい場合にしか結果の変わらない判定、どのモデルも呼ばないコード）。したがって検出率は83個中74個（\SI{89}{\percent}）、サニタイザを含めると76個（\SI{92}{\percent}）である。残る7個の内訳は次のとおりである。2個は結果を丸め誤差の範囲か誤差の定数でしか変えない（\val{"CZ"}の壁での面内拡散、$1/\mu$の深さ平均におけるゴースト層）。2個はクーラン数の警告しか変えない。3個はどのテストも到達しない振る舞いであった。
  \item 入出力コードからの25個のうち、\code{ctest}は20個を検出した。2個は何も変えようがない（何からも呼ばれないコード）ので、23個中20個（\SI{87}{\percent}）である。サニタイザで増えた分はない。残る3個は、\val{"CZ"}の\file{.vts}の最外ゴースト層の速度、格子サイズの上限、ファイル末尾まで続く\file{.inp}文字列の読み取りであった。2つの層の大きさで重み付けすると\SI{89}{\percent}、サニタイザを含めると\SI{91}{\percent}である。
\end{itemize}

その後、次の確認を追加した。濃度にKorenスキームを用いた1次元と2次元の鉛直一様モデルの実行の比較と、細長いセル上で転置したノズル（\code{equivalence}）、場ファイルから読んだ片側だけ$[0, 1]$の外にある濃度、格子サイズの上限、終端のない\file{.inp}文字列（\code{input\_guards}）、2次元ノズルの面積（\code{analytic}）である。これらは生き残ったもののうち四個を検出し、標本の検出数は83個中77個（サニタイザを含めて79個）と23個中21個になった。この過程で欠陥が二つ見つかり、どちらも修正した。一つは2次元ノズルで、幅の3.9倍より長いセルの上では部分的に覆うセルを除外しており、幅の11倍長いセルでは流量の\SI{81}{\percent}しか供給していなかった（\cref{sec:verification:analytic}）。もう一つは閉じ引用符のない\file{.inp}文字列で、ファイルの残り全体を値として取り込んでいた。後者の発見につながった変更は、修正後の読み取り処理では何も変えない。

第四の標本は、第三の標本の後に変更したコードから抽出した。対象は、落ち着かない\val{"CZ"}の結合パス、$c = 1$を超えた分から$c = 0$を下回るセルへ移す溶質、保存形の鉛直一様モデルの溶質収支とその開いた辺、停滞するニュートン修正、64ビットのステップ数である。第三の標本以降に変更した行にある511通りの変更から、以前に抽出したものを除いて30個（シード20261006）を取り、それぞれについて\code{ctest}全体を実行した。

\begin{itemize}
  \item \code{ctest}は16個を検出し、うち1個は制限時間の超過によるものであった。十個は結果を変えようがない。そのうち六個は、厳密に等しい場合か、どの入力も到達しない場合にしか違いが出ない。残る四個は、ループや添字を誰も読まないゴーストセルやスロットへ動かすもので、サニタイザも報告しない。したがって検出率は20個中16個（\SI{80}{\percent}）、サニタイザを含めると17個（\SI{85}{\percent}）である。サニタイザは、$c = 0$を下回る\val{"CZ"}の列を埋める処理で、末尾を越えて回るループを検出する。
  \item 残る三個のうち二個は、どのテストも到達しない振る舞いであった。2次元でその列の隣の列が$y$方向に逆側から読む変更と、すべてのステップを溶質を移したステップとして数える変更である。三個目は、結合パスのデバッグ出力しか変えない。
\end{itemize}

その後、壁際で$c = 0$を下回る状態からリスタートした\val{"CZ"}の列を、その転置と比較する確認（\code{equivalence}の\code{cz2d\_deficit\_transpose}）を追加した。これは両方を検出し、標本の検出数は20個中19個となった。欠陥は見つからなかった。

\section{性能と規模}
\label{sec:notverified:scale}

\begin{itemize}
  \item \textbf{実時間と並列スケーリングは\cref{ch:verification}では測定していない。}同章のバックエンドの比較は、値と反復回数の一致を見るだけである。\footnote{\cref{tab:linsolve-timing}の求解時間の目安は、バックエンドを選ぶ参考として\cref{sec:linsolve:choosing}に載せたものである。検証のための測定ではなく、\cref{ch:verification}の再現性の保証はいずれも適用されない。}
  \item \textbf{検証に用いた格子は、最大で1次元が$N_x = 4000$、2次元が$N_x = N_y = 120$である。}より大きな格子（$N_x = 10^{6}$、$1024^{2}$、$200^{2} \times 20$）は、メモリ使用量を測るために数ステップだけ実行した（\cref{sec:verification:long}）。これらの格子で長時間実行したときの振る舞い（線形ソルバーの収束、丸め誤差の蓄積）は測定していない。
\end{itemize}
